% 第 12 题：两挡位欧姆表的内部结构（TikZ 电路）
% G 与 R0 上段串联为上支路，Rs 与 R0 下段串联为下支路，两支路并联后经 Rm 或 Rn 接开关 S
\begin{tikzpicture}[font=\small, >=Stealth]
  % 电池与端子①（电池接在并联组合的下端）
  \draw (0,4.4) -- (0,2.1);
  \draw (0,2.1) to[battery1] (0,0.6);
  \draw (0,0.6) -- (0,0.35);
  \fill (0,0.35) circle (1.6pt);
  \node[below] at (0,0.3) {①};
  \fill (0,4.4) circle (1.8pt);
  % 上支路：G（导线止于 G 的两侧）
  \draw (0,4.4) -- (0.61,4.4);
  \draw (0.95,4.4) circle (0.34) node {$\mathrm{G}$};
  \draw (1.29,4.4) -- (4.15,4.4) -- (4.15,4.1);
  % 调零电阻 R0
  \draw (4.15,4.1) rectangle (4.70,2.10);
  \node[left] at (4.05,3.1) {$R_0$};
  \draw[->] (5.35,3.1) -- (4.72,3.1);
  \fill (5.35,3.1) circle (1.8pt);
  % 下支路：Rs
  \draw (0,2.10) -- (1.00,2.10);
  \draw (1.00,1.85) rectangle (2.60,2.35);
  \node[below] at (1.80,1.85) {$R_s$};
  \draw (2.60,2.10) -- (4.15,2.10);
  % 右侧节点
  \draw (5.35,3.1) -- (6.95,3.1) -- (6.95,2.85);
  \draw (5.35,3.1) -- (5.35,2.85);
  % Rm
  \draw (5.10,2.85) rectangle (5.60,1.65);
  \node[left] at (5.00,2.25) {$R_m$};
  \draw (5.35,1.65) -- (5.35,1.05);
  \fill (5.35,0.88) circle (1.5pt);
  \node[below] at (5.35,0.82) {$m$};
  % Rn
  \draw (6.70,2.85) rectangle (7.20,1.65);
  \node[right] at (7.30,2.25) {$R_n$};
  \draw (6.95,1.65) -- (6.95,1.05);
  \fill (6.95,0.88) circle (1.5pt);
  \node[below] at (6.95,0.82) {$n$};
  % 开关 S 与端子②
  \draw (6.60,-0.15) -- (5.55,0.55);
  \fill (6.60,-0.15) circle (1.5pt);
  \node[right] at (6.70,-0.15) {$\mathrm{S}$};
  \draw (6.60,-0.15) -- (6.60,-0.65);
  \fill (6.60,-0.85) circle (1.6pt);
  \node[below] at (6.60,-0.9) {②};
\end{tikzpicture}
